\chapter{Adrenal 21-Hydroxylase Antibodies}

\section*{What Is This Test?}
Adrenal 21-hydroxylase antibodies are autoantibodies directed against a key enzyme in adrenal steroid synthesis. They are the principal serologic marker of autoimmune adrenalitis, the most common cause of primary adrenal insufficiency in many populations.

\section*{Alternative Names}
21-Hydroxylase Antibody; CYP21A2 Antibody; Adrenal Cortex Antibody.

\section*{Why This Test Is Ordered}
Testing helps determine whether confirmed primary adrenal insufficiency is autoimmune and can identify people at increased risk in selected autoimmune-polyglandular settings. It does not replace cortisol, ACTH, or stimulation testing when adrenal insufficiency is suspected.

\section*{How This Test Is Performed}
A blood sample is analyzed by an immunoassay for antibodies against 21-hydroxylase. Results may be interpreted with cortisol, ACTH, renin, aldosterone, electrolytes, thyroid antibodies, and other autoimmune evaluation.

\section*{Preparation, Risks, and Limitations}
No fasting is required. Venipuncture is the usual risk. Antibody levels can decline over time, and a negative result does not exclude nonautoimmune primary adrenal disease or every autoimmune case.

\section*{Understanding Results}
A positive result supports autoimmune adrenalitis in a compatible clinical setting and may prompt assessment for associated autoimmune disease. A negative result shifts attention toward genetic, infectious, infiltrative, medication-related, or other causes.

\section*{References}
Endocrine Society guidance on diagnosis and treatment of primary adrenal insufficiency.

\clearpage
\section*{Understanding Results and Follow-up}
The laboratory report should identify the specimen, method, units, reference interval or decision limit, and whether the result is positive, negative, abnormal, or inconclusive. Reference intervals vary with age, sex, pregnancy, specimen type, assay, and laboratory; a result outside a range is not automatically a diagnosis, and a result within a range does not always exclude disease. Discuss unexpected results with the ordering clinician, who can decide whether repeat or confirmatory testing, treatment, or referral is appropriate. Do not change medicines or delay urgent care based on an isolated result.


\section*{Regulatory Status and Authorization Date}
Regulatory status applies to the specific assay, instrument, manufacturer, and intended use rather than to the test name alone. No single approval date can be assigned to this test category. For a commercial assay, verify the current product in the relevant regulator's database: the U.S. Food and Drug Administration (FDA) clearance, approval, or authorization; India's Central Drugs Standard Control Organisation (CDSCO) licence or permission; the European Union In Vitro Diagnostic Regulation (EU IVDR) CE certificate issued through the applicable conformity-assessment route; or the United Kingdom Medicines and Healthcare products Regulatory Agency (MHRA) registration and UKCA/CE status. Laboratory-developed tests may not have an individual FDA, CDSCO, EU IVDR, or MHRA approval date.
\clearpage
